\documentclass[11pt,a4paper]{article}
\usepackage[T1]{fontenc}
\usepackage[utf8]{inputenc}
\usepackage[margin=2.3cm,top=2.6cm,bottom=2.6cm]{geometry}
\usepackage{lmodern}
\usepackage{microtype}
\usepackage{xcolor}
\usepackage{booktabs}
\usepackage{enumitem}
\usepackage{titlesec}
\usepackage{fancyhdr}
\usepackage{tcolorbox}
\usepackage[hidelinks]{hyperref}

\definecolor{ink}{HTML}{16201B}
\definecolor{accent}{HTML}{1A6B3C}
\definecolor{muted}{HTML}{5D6B64}
\definecolor{rule}{HTML}{D7DED9}
\definecolor{soft}{HTML}{F4F7F5}
\definecolor{warn}{HTML}{9A6B00}

\color{ink}
\titleformat{\section}{\Large\bfseries\color{accent}}{\thesection.}{0.5em}{}
\titleformat{\subsection}{\large\bfseries}{\thesubsection}{0.5em}{}
\titlespacing*{\section}{0pt}{18pt}{7pt}
\titlespacing*{\subsection}{0pt}{13pt}{4pt}

\pagestyle{fancy}\fancyhf{}
\renewcommand{\headrulewidth}{0.4pt}
\fancyhead[L]{\small\color{muted}ReqShift and GoalShift}
\fancyhead[R]{\small\color{muted}\thepage}

\setlist[itemize]{leftmargin=1.3em,itemsep=2pt,topsep=3pt}
\setlist[enumerate]{leftmargin=1.5em,itemsep=2pt,topsep=3pt}

\newtcolorbox{keybox}[1][]{colback=soft,colframe=accent,boxrule=0.7pt,
  arc=2pt,left=8pt,right=8pt,top=6pt,bottom=6pt,#1}
\newtcolorbox{warnbox}[1][]{colback=white,colframe=warn,boxrule=0.7pt,
  arc=2pt,left=8pt,right=8pt,top=6pt,bottom=6pt,#1}

\newcommand{\gs}{\textbf{GoalShift}}
\newcommand{\rs}{\textbf{ReqShift}}

\begin{document}

\thispagestyle{empty}
\begin{center}
\vspace*{2.2cm}
{\Huge\bfseries\color{accent} ReqShift and GoalShift}\\[10pt]
{\LARGE What they are, how they work,\\[3pt] and how they might become a product}\\[26pt]
{\large Daniyal Bukhari}\\[6pt]
{\color{muted}\today}
\end{center}

\vspace{1.4cm}

\begin{keybox}
\textbf{The one sentence version.} A language model cannot tell you when it is wrong.
A SAT solver cannot read prose. Their weaknesses are opposites, so \rs{} uses AI to turn
a requirements document into a formal model, and \gs{} then proves what follows from it.
Nothing the tool \emph{claims} comes from the AI.
\end{keybox}

\vspace{0.6cm}
{\color{muted}\small
This document has two halves. Sections 1 to 5 describe the software as it stands today,
with real measured numbers. Sections 6 to 9 are about turning it into something
sellable, and those are ideas and options rather than facts. Where something is
uncertain or unverified, it says so.}

\newpage
\tableofcontents
\newpage

\section{The short version}

You have two pieces of software. They do different jobs and it matters that they stay
separate.

\begin{center}
\begin{tabular}{@{}p{4.2cm}p{5.1cm}p{5.1cm}@{}}
\toprule
 & \textbf{GoalShift} & \textbf{ReqShift} \\
\midrule
What you give it & A goal model you drew & A requirements document \\
What it gives back & Behavior classes, a statechart & A goal model, plus a report \\
Who does the thinking & A SAT solver & AI, then GoalShift \\
Can it be wrong? & No. Within your model it is proved & Yes. The extracted model is a draft \\
Needs internet & No & Yes \\
In your thesis and paper & \textbf{Yes} & \textbf{No} \\
Size & 445 lines of Python, 1{,}923 lines of HTML & about 3{,}300 lines across 11 files \\
\bottomrule
\end{tabular}
\end{center}

\vspace{6pt}

Think of a maths word problem. \gs{} is the calculator: give it the equation and it gives
the exact answer, but it cannot read the word problem. \rs{} is the part that reads the
word problem and writes down the equation, but it cannot do the arithmetic.

\begin{warnbox}
\textbf{Why they must stay separate.} \gs{} is the artifact your thesis describes and
your paper will describe. Its SAT encoding \emph{is} the contribution. If you change it,
the numbers in your own chapter stop being reproducible. \rs{} imports \gs{} read only
and never modifies it. Delete \rs{} tomorrow and \gs{} is exactly as it was.
\end{warnbox}

\section{How GoalShift works}

\subsection{What goes in}

A \emph{restricted type-level goal model}. In plain terms, a description of who is in the
system, what each one is trying to achieve, and who needs help from whom.

\begin{itemize}
\item \textbf{Actor types} with a cardinality range. Not individual actors. ``Up to five
      ambulances'' is one actor type with an upper bound of five.
\item \textbf{Goals} belonging to exactly one actor, one of them marked as that actor's
      root goal.
\item \textbf{Decomposition} breaking a goal into sub-goals of the same actor, using
      AND, inclusive OR, or XOR.
\item \textbf{Dependencies}, the only link allowed between different actors. A
      dependency carries a multiplicity: \texttt{1..*} means at least one present
      provider must supply it, \texttt{*..*} means every present provider must.
\end{itemize}

\subsection{The four steps}

\begin{enumerate}
\item \textbf{Encode to CNF.} Every actor instance and every goal instance becomes a
      Boolean variable. Four groups of clauses are generated: cardinality (including
      symmetry breaking so present instances form a prefix), decomposition, dependency,
      and linking (an absent actor satisfies nothing). This formula is $\Phi_{TGM}$.
\item \textbf{Walk the configurations.} A configuration is how many of each actor type
      are present. For each one, small SAT checks read off which actors can meet their
      root goal.
\item \textbf{Group into behavior classes.} Configurations with the same satisfaction
      pattern collapse into one class. These are equivalence classes, so testing one
      member tests them all.
\item \textbf{Derive the statechart.} Classes become states, and actors joining or
      leaving become transitions.
\end{enumerate}

\subsection{The scalability fix}

This is worth recording because it changed what the tool can do.

The original engine answered ``which behavior classes exist'' by asking the solver to
\textbf{list every complete satisfying assignment} and then grouping them. For the small
running example that was 206 assignments, used to report 4 classes. The count roughly
doubles with every added goal or actor, so it went from instant to never finishing with
no warning in between.

\begin{center}
\begin{tabular}{@{}lrr@{}}
\toprule
\textbf{Dimension} & \textbf{Old ceiling} & \textbf{New ceiling} \\
\midrule
Actor types & 3 & 10 (7.5 s) \\
Instance actors per type & 2 & 30 (7.6 s) \\
Goals per actor & 3 & 8{,}191, at 49{,}152 variables (0.6 s) \\
\bottomrule
\end{tabular}
\end{center}

The replacement is a \emph{projected} query. Only the actor-presence variables and the
actor root goals are ever read, so the solver is asked about those directly and left to
decide the rest internally. Goal tree size stops mattering, because a solver decides
\emph{whether} a big tree can be satisfied about as fast as a small one.

\begin{keybox}
\textbf{The point for a defense or a paper.} The approach in the thesis was never the
problem. The prototype's way of querying the solver was wasteful. Verified by running
the old and new engines side by side on 29 generated models and comparing the full
result, instance table, classes and statechart: \textbf{zero mismatches}.
\end{keybox}

\subsection{The reference numbers}

The Chapter 7 emergency coordination case study, which is the regression target:

\begin{center}
\begin{tabular}{@{}lrrrrr@{}}
\toprule
\textbf{Engine} & \textbf{Variables} & \textbf{Clauses} & \textbf{Configurations}
  & \textbf{Solver calls} & \textbf{Classes} \\
\midrule
GoalShift (SAT) & 47 & 170 & 24 & 86 & 4 \\
z3 (SMT) & 15 & 22 & 24 & 12 & 4 \\
\bottomrule
\end{tabular}
\end{center}

\section{How ReqShift works}

\subsection{The pipeline}

\begin{center}
\begin{tabular}{@{}ll@{}}
\toprule
\textbf{Stage} & \textbf{Who does it, and can it be wrong} \\
\midrule
Requirements document & your input \\
$\downarrow$ AI extraction & \textbf{a draft.} Review it \\
Goal model (GoalShift's own JSON) & \\
$\downarrow$ human review on the canvas & two minutes of your time \\
$\downarrow$ GoalShift and SAT & \textbf{proved} \\
Behavior classes, statechart, findings & \\
$\downarrow$ AI write-up (optional) & rewording of proved facts only \\
Report & \\
\bottomrule
\end{tabular}
\end{center}

The AI appears only at the two ends. The middle is your thesis and is untouched.

\subsection{The files}

\begin{center}
\begin{tabular}{@{}lrp{8.4cm}@{}}
\toprule
\textbf{File} & \textbf{Lines} & \textbf{Job} \\
\midrule
\texttt{schema.py} & 224 & The JSON schema the AI must fill in, plus conversion to GoalShift's format \\
\texttt{extractor.py} & 263 & The only file that talks to a model. Two backends, one schema \\
\texttt{analyze.py} & 173 & Imports GoalShift, derives findings. \textbf{No AI involved} \\
\texttt{analyze\_smt.py} & 632 & The experimental z3 backend. Opt in \\
\texttt{report.py} & 213 & Markdown report and the plain-text brief. \textbf{No AI involved} \\
\texttt{reqshift.py} & 197 & The command line tool \\
\texttt{app.py} & 282 & The web interface's server \\
\texttt{ui/index.html} & 964 & The whole interface, one self-contained file \\
\texttt{viewer.py} & 144 & Read-only browser for reports already on disk \\
\bottomrule
\end{tabular}
\end{center}

\subsection{Why the AI cannot invent a defect}

Every finding is computed arithmetically from the solver's output, inside
\texttt{analyze.py}, which never calls a model at all. The optional write-up pass is
given only the already-proved facts and told it may reword them but not add to them.

This is the whole design, and it is the answer to ``how do I know your AI did not
hallucinate this?''

\subsection{What it finds}

Four findings from the SAT path:

\begin{itemize}
\item \textbf{blocked}: an actor can be present and provably unable to meet its root
      goal. The expensive one, because the equipment is switched on and doing nothing.
\item \textbf{never satisfied}: an actor reaches its root goal in no configuration at
      all. Usually an extraction mistake or an unachievable requirement.
\item \textbf{initial degraded}: the minimum configuration, the state the system boots
      into, does not satisfy everything present.
\item \textbf{inert cardinality}: a range the requirements allow that changes no
      outcome anywhere.
\end{itemize}

Two more become possible with the SMT path, and both are \emph{solved for} rather than
observed, so the numbers are exact:

\begin{itemize}
\item \textbf{capacity limit}: ``Treatment Van is capped at 4 by Clinic Site capacity,
      not by its own bound of 8.'' The declared upper bound is not the real one.
\item \textbf{deadline miss}: ``cannot meet its 60 minute deadline, the soonest it can
      finish is 85.''
\end{itemize}

\subsection{Two ways to reach the AI}

\begin{itemize}
\item \textbf{CLI backend}: uses the \texttt{claude} command from Claude Code. Needs no
      API key if you are already signed in. This is the zero-setup path and the one used
      for every result in this document.
\item \textbf{SDK backend}: the official Anthropic SDK with structured outputs, which
      constrains the reply to the schema at the API level rather than by asking nicely.
      Needs an API key. Written from the documentation but \textbf{not yet tested},
      because there is no key on this machine.
\end{itemize}

\section{The SMT extension}

\subsection{Integer counts}

\texttt{analyze\_smt.py} analyses the same models with z3 integer variables instead of
one Boolean per possible instance. \gs{} stays the default and stays untouched.

Two things follow. First, the lower bound may exceed one, so ``at least three robots''
becomes expressible where the SAT encoding rejects it outright. Second, and more
interesting, a structural fact about the notation emerges:

\begin{keybox}
Every place a count is read, the only question asked is whether it is zero.
A \texttt{1..*} dependency needs $n \geq 1$, a \texttt{*..*} dependency is vacuous at
$n = 0$, and goals are false at $n = 0$. \textbf{Nothing in the notation can tell two
present hospitals from three.} So the satisfaction pattern depends on the zero/nonzero
vector alone, of which there are at most $2^k$ for $k$ actor types, however wide the
ranges are.
\end{keybox}

That also \emph{explains} the inert cardinality finding rather than merely detecting it,
and it means the classes come back symbolic, which is a better answer than a list:

\begin{center}\ttfamily\small
Class 4 (OK, OK, OK), 15 configurations\\
holds when n\_D = 1 and 1 <= n\_H <= 3 and 1 <= n\_M <= 5
\end{center}

Four lines say everything a 24 cell grid says, and they stay four lines if the bounds go
to a million. Measured: one million maximum instances per type gives $10^{18}$
configurations in 12 solver calls.

\subsection{Timing and capacity}

Four optional numeric fields, filled by the extractor only when the text states the
figure: \texttt{duration} on a goal, \texttt{deadline} on an actor, \texttt{capacity} on
an actor, \texttt{demand} on a dependency.

Adding any of them changes the shape of the problem. Once a real number is in play the
answer stops depending only on whether a count is zero, so the fast path is dropped and
configurations are walked. That is the honest cost.

What it buys is real. On the case study with each hospital admitting one patient, a
\textbf{fifth behavior class appears} that cannot exist in the Boolean model: hospitals
fine, ambulances present, ambulances blocked because there are not enough beds.

\subsection{How it is kept honest}

The SMT encoding rests on one assumption, that present instances of an actor type can be
taken to behave alike when deciding that type's satisfaction. That should be exact here,
since instances are interchangeable and share no constraint, but it is not proved.

So \texttt{bench/diff\_smt.py} runs \gs{} as the oracle: \textbf{77 models compared, zero
mismatches}, including the hand-built Chapter 7 model, four models extracted from real
text, and 72 generated ones.

\begin{warnbox}
A model analysed with the SMT engine is \textbf{no longer proved by the encoding the
thesis describes}. The interface shows an amber banner saying so. Keep that label on
anything you show anyone.
\end{warnbox}

\section{Evidence that it works}

\subsection{It finds real defects}

\texttt{examples/welding\_cell.txt} is a plausible robotic welding cell specification.
Read it and nothing looks wrong.

From 349 words the extractor produced a model, and the solver proved that in
\textbf{8 of 20 configurations} a piece of equipment is present and cannot do its job,
including the configuration the cell boots into. The cause, traced back to the
sentences responsible: two requirements set \emph{upper} bounds on robots and vehicles
and never a lower one, so zero is legal and everything downstream stalls.

It also flagged that the numbers four and three carry no behavioural consequence at all,
which is worth knowing before someone plans capacity around them.

\subsection{It reproduces a model you built by hand}

Your Chapter 7 case study, from the 143 words of section 7.2 alone, which contains no
goal model, no AND/OR/XOR, and no multiplicity notation.

Three independent runs, one from the command line and two through the interface, all
produced the same four behavior classes with the same 1/3/5/15 membership split as your
hand-built model. Compared by actor role rather than column position, the result is
\textbf{identical}.

\begin{keybox}
\textbf{Worth putting in the paper.} The extraction consistently placed the
\texttt{*..*} dependency on a different edge than you did, and invented a dispatch-side
goal because your section 7.2 describes hospitals reporting to the dispatch center
without giving the center a goal that needs it. Two structurally different type-level
models, identical behavior classes, reproduced across runs. That is evidence the
behavior-class abstraction is robust to how a modeller draws the dependency.
\end{keybox}

\subsection{It reads numbers out of prose}

\texttt{examples/dialysis\_unit.txt} states its figures in ordinary sentences. From 314
words the extractor pulled all six durations converted to one unit, a 120 minute
deadline, capacity four from ``each clinic site has four treatment bays'', and demand one
from ``each van occupies one bay''.

The timing was checked by hand: $6 + 4 + 22 + 45 + 8 = 85$ minutes, which is exactly what
z3 returned as the soonest possible finish. The latency chain composes correctly across
the dependency.

\newpage
\section{From prototype to product: the honest starting position}

Everything above is real and measured. What follows is opinion and options.

\subsection{What you have}

\begin{itemize}
\item A working analyzer with a validated formal foundation and a thesis behind it.
\item An AI front end that removes the step that stops people using formal tools, namely
      building the model by hand.
\item A differentiator that is genuinely uncommon. Plenty of tools do ``AI reads
      requirements''. Very few reason about \emph{membership change}, and that is your
      actual contribution.
\item Evidence, not claims. Differential tests, reproduced results, measured ceilings.
\end{itemize}

\subsection{What you do not have}

Being clear about this is more useful than optimism.

\begin{itemize}
\item \textbf{No product.} A prototype that a knowledgeable author can run is a long way
      from software a stranger can rely on.
\item \textbf{No customer.} Nobody has paid, and nobody has used it on their own
      document without you present.
\item \textbf{No company, and no clarity on who owns the work.} See below.
\item \textbf{Real notation limits.} No soft goals, no contribution links, lower bounds
      restricted in the SAT path, XOR restricted to two children, and sequencing lost
      to AND.
\item \textbf{The extraction is a draft.} This is by design and it is the right design,
      but it means the tool assists an expert rather than replacing one.
\end{itemize}

\begin{warnbox}
\textbf{Sort the intellectual property before anything else.} The work grew out of a
thesis at TH Bingen, supervised by Prof.\ Bandyszak, and a paper is being written jointly.
Rules on student and university IP vary by institution and by how the work was funded.
Ask the university's technology transfer office in writing, and agree with your advisor
what is shared and what is yours, \emph{before} you show it to a company. This is not
legal advice and you should get some.
\end{warnbox}

\section{Four routes, roughly in order of how well they fit you}

\subsection{Route 1: a German founder grant (best fit)}

The \textbf{EXIST-Gr\"understipendium} is a federal programme for founders coming out of
universities. It typically provides a monthly stipend for graduates for twelve months,
plus money for materials and coaching, and requires a university to host the project.

Why it fits you specifically: a fresh master's thesis with a commercialisable result, an
advisor who could host it, and a technology that came out of academic work. That is
close to the profile the programme exists for.

What to do: talk to the TH Bingen transfer office and to Prof.\ Bandyszak about hosting.
\textbf{Verify the current terms yourself}, since amounts and rules change and I cannot
confirm today's figures. There is also \textbf{EXIST-Forschungstransfer} for
research-heavy projects, which is larger and harder to get.

\subsection{Route 2: sell the analysis before selling the tool}

The fastest way to learn whether anyone wants this is to do the work by hand for money.

Offer a fixed price review: they send one real specification, you run the tool, you
review the extracted model yourself, and you hand back a report naming the defects and
the sentences responsible. Perhaps five to fifteen thousand euro for one specification.

This is far easier to sell than software, because it is a service with a clear
deliverable and no integration. It also tells you exactly what to build, since you will
discover on the first document what the notation cannot express.

\subsection{Route 3: design partners, not customers}

Before charging anyone, find three companies willing to give you a real specification
and two hours of an engineer's time, in exchange for a free analysis.

You are not selling. You are finding out whether the defects you surface are ones they
care about. If an engineer says ``we already knew that'', the wedge is wrong. If they go
quiet and then ask who else has seen it, the wedge is right.

\subsection{Route 4: build it inside a company}

Joining a company already selling into this space and building it there removes the
funding problem entirely, at the cost of ownership.

BoWatt is the obvious example and they rejected you, but for a Full Stack Engineer role,
not for this. That is worth separating. A message that says ``you build agentic
requirements engineering, I wrote a thesis on formally verifying multi-actor requirements
models and here is a report my tool produced from a spec'' is a completely different
approach from a job application, and it can be sent to the founders directly.

\section{Who might buy this, and what to say}

\subsection{The positioning that is not crowded}

Do not say ``AI for requirements''. That market is loud and you will be compared to
better-funded companies.

Say instead:

\begin{keybox}
\emph{We prove which configurations of your system still work when parts of it are
missing or unavailable, and we cite the requirement sentence that causes each failure.}
\end{keybox}

Two claims in there are hard for competitors to match. \emph{Prove} is real, because a
solver produced it. \emph{Cite the sentence} is real, because traceability is built in.

\subsection{Why certification markets are the right target}

Cars, planes, medical devices and industrial safety all live under standards that demand
traceability and evidence: ISO 26262, DO-178C, IEC 61508. In those markets, formal
evidence is not a nice-to-have, it is something they already pay people to produce.

Your test reduction argument lands here too. Behavior classes are equivalence classes,
so a handful of scenarios provably cover every configuration. A validation engineer
understands the value of that immediately.

\subsection{Where these people actually are}

\begin{itemize}
\item \textbf{Fraunhofer IESE, Kaiserslautern.} Does requirements engineering research
      directly with industry, in Germany, near you. Probably the single most useful
      first conversation available to you, because they know both the research and the
      companies.
\item \textbf{Conferences with industry tracks:} REFSQ (requirements engineering,
      European, industry-friendly), IEEE RE, MODELS, INCOSE symposiums, and ERTS for
      embedded and real-time systems.
\item \textbf{Tool vendors as partners rather than competitors:} IBM DOORS, Siemens
      Polarion, PTC codebeamer, Jama Connect. These own the workflow already. Being an
      analysis that plugs into their ecosystem is a much shorter path than replacing
      them.
\item \textbf{Your advisor's network.} Prof.\ Bandyszak's own MODELSWARD paper on
      type and instance goal models came from somewhere, and that somewhere has
      industrial contacts.
\end{itemize}

\subsection{What BoWatt appears to have done}

From public information only. They bootstrapped to over half a million euro of annual
recurring revenue \emph{before} raising their first round, and only then took money from
tier-one investors and angels.

The lesson worth copying is the order. They sold something to somebody first. Revenue
before funding is also the version of this that does not require you to convince an
investor of anything.

\section{Roadmap, if you decide to pursue it}

Ordered so that each step is useful even if you stop after it.

\begin{enumerate}
\item \textbf{Publish the paper with Torsten.} Credibility first, and it costs you
      nothing extra since it is already in progress.
\item \textbf{Settle the IP question.} In writing.
\item \textbf{Run it on three real specifications} from three different companies, free.
      Learn what breaks.
\item \textbf{Fix the notation gaps those documents expose.} Almost certainly soft goals
      and better handling of ordering.
\item \textbf{Make the extraction reviewable by someone who is not you.} The canvas
      review step is the product, and right now it needs a console paste.
\item \textbf{Charge for one analysis.} The first invoice tells you more than any amount
      of planning.
\item \textbf{Only then} consider whether it is a company, a grant application, a
      licence to a tool vendor, or a job at a company that already sells here.
\end{enumerate}

\section{Risks worth saying out loud}

\begin{itemize}
\item \textbf{Enterprise sales cycles in safety-critical industries run six to eighteen
      months.} Plan funding accordingly, which is another argument for the grant route.
\item \textbf{Domain credibility matters more than technology} in these markets, and a
      recent graduate has little of it. A partner, an advisor, or an institute
      affiliation compensates.
\item \textbf{The incumbents own the workflow.} Nobody rips out DOORS. Integration beats
      replacement.
\item \textbf{Your notation restrictions are real} and the first industrial document you
      try will expose several at once. That is not failure, it is the roadmap.
\item \textbf{The tool cannot tell a defect from a design decision.} In the welding cell,
      a degraded boot configuration was a real bug. In your case study, the identical
      finding is the intended result. Same output, opposite meaning. Say that before
      anyone else does.
\end{itemize}

\vspace{10pt}
\begin{keybox}
\textbf{The honest summary.} You built the hard half. Anyone can bolt an AI onto the
front of a tool. Almost nobody has the thing behind it that can prove the answer. What
you do not yet have is a customer, and everything above is really just different ways of
going and finding one.
\end{keybox}

\end{document}
